\documentclass[11pt,a4paper]{article}

\usepackage[utf8]{inputenc}
\usepackage[T1]{fontenc}
\usepackage[brazil]{babel}
\usepackage[margin=2.4cm]{geometry}
\usepackage{amsmath,amssymb}
\usepackage{booktabs}
\usepackage{xcolor}
\usepackage{graphicx}
\usepackage[hidelinks]{hyperref}

\newcommand{\code}[1]{\texttt{\small #1}}
\newcommand{\unit}[2]{\ensuremath{#1\,\mathrm{#2}}}

\title{\textbf{Por que os dados de campo do GRAVI-SH não se encaixam
direto nas normas}\\[0.3em]
\large Causas da divergência entre o kit de aquisição do laboratório e a
ISO 2631-1/2631-5/NHO-09, e mudanças possíveis}
\author{GRAVI-SH --- Grupo de Acústica e Vibrações em Seres Humanos\\
Universidade Federal de Minas Gerais}
\date{\today}

\begin{document}
\maketitle

\begin{abstract}
\noindent
Este apontamento lista os motivos pelos quais os dados do kit de aquisição
de campo do GRAVI-SH --- descrito em Goeking (2023), Goeking et al. (2024)
e Avila et al. (2026) --- não se encaixam nas cadeias de cálculo
normativas (Anexo A/E da ISO 2631-5, NHO-09, Diretiva 2002/44/CE), e o que
isso implica para o \emph{Avaliador de Vibração Humana}. As causas estão
no dado de origem, não no software. O caminho de campo do software
(Cláusula 5/NHO-09/Diretiva, sempre não-normativo fora do assento)
continua sendo a escolha correta hoje.
\end{abstract}

\section{Contexto}

O \emph{Avaliador de Vibração Humana} processa dois tipos de dado. O
primeiro é o registro ASCII de exemplo da norma ou de bancada controlada:
um arquivo contínuo, com os quatro pontos de contato (assento, encosto,
pés, mãos) amostrados de forma simultânea, como o Anexo A pressupõe. O
segundo é o arquivo \code{.m} do registrador de campo do GRAVI-SH, montado
no usuário por um dia inteiro fora de laboratório. Este apontamento trata
do segundo tipo.

O kit está descrito em Avila et al. (2026, CONEM2026-0596), que dá
sequência a Goeking (2023) e Goeking et al. (2024): três unidades
independentes (assento, encosto, apoio de pé), cada uma com acelerômetro
MPU6050, microcontrolador ESP8266 (placa WeMos Mini D1), relógio de tempo
real DS1302 e cartão microSD, amostrando a \unit{1024}{Hz}, amostras de 16
bits. É o formato lido por \code{iso26315/fielddata.py}.

\section{Motivos dos problemas}

\subsection{Ciclo de trabalho de 40\%}
\label{sec:duty}

O firmware grava blocos de \unit{2}{s} a cada \unit{5}{s} de relógio. Os
\unit{3}{s} restantes são usados para escrever no cartão SD os dados
aferidos nesses \unit{2}{s} -- a escrita não é instantânea, e essa é a
causa do ciclo de 40\%, não uma escolha de economia de energia ou memória.

\paragraph{Por que isso quebra a dose do Anexo A.} O Anexo A calcula a
força de compressão dinâmica em cada nível discal aplicando a função de
transferência do modelo de elementos finitos (Fórmula 1 do Anexo A) aos
quatro pontos de contato (assento, encosto, pés, mãos), simultâneos. Do
sinal de força resultante, extraem-se os picos de compressão entre
cruzamentos consecutivos da média ($C_{dyn,i}$, A.3.3), e a dose por
exposição é

\begin{equation}
S = \left(\sum_i \left(\frac{|C_{dyn,i}|}{B}\right)^{6}\right)^{1/6}
\qquad \text{(Fórmula A.1)}
\end{equation}

com $B$ a área da placa terminal vertebral. A dose diária equivalente
combina $j$ exposições medidas:

\begin{equation}
S_d^A = \left(\sum_j S_j^{6}\,\frac{t_{dj}}{t_{mj}}\right)^{1/6}
\qquad \text{(Fórmula A.3)}
\end{equation}

O risco $R^A$ (Fórmulas E.1/E.4) tem a mesma estrutura: soma de sexta
potência ao longo dos anos de exposição, dividida pela diferença entre a
resistência óssea $S_u$ e a tensão estática.

O ponto central: $S$ soma sobre os picos $C_{dyn,i}$ \emph{que existem no
sinal}. Se 60\% do tempo não foi gravado, os picos que ocorreriam nesse
intervalo simplesmente não entram na soma -- não é um problema de ruído ou
de precisão, é ausência do dado que a fórmula soma. Como o expoente é 6, a
dose é dominada pelos poucos maiores picos, não por uma média; a
extrapolação $t_d/t_m$ da Fórmula A.3 assume que o trecho medido
representa o padrão de picos do período não medido, premissa razoável
para uma medição contínua de alguns minutos representando um turno de
trabalho, mas não para um sinal que já nasce com 60\% do tempo ausente de
forma estrutural. Por isso o software não aplica o Anexo A/E a dados de
campo.

\subsection{A Cláusula 5 (Anexos C e D) tem o mesmo problema}
\label{sec:clause5}

A Cláusula 5 evita o problema de sincronismo da
Seção~\ref{sec:sync} porque usa só o eixo z do assento, com uma função de
transferência fixa (Fórmula 1 da Cláusula 5, Anexo D dá a implementação em
filtro digital causal dessa mesma função) -- não precisa combinar quatro
pontos simultâneos. Mas a estrutura da dose é a mesma do Anexo A:

\begin{equation}
D_z = 1{,}07\left(\sum_i A_i^{6}\right)^{1/6}
\qquad \text{(Fórmula 3)}
\end{equation}

soma de sexta potência dos picos positivos $A_i$ da resposta da coluna, e

\begin{equation}
D_{zd} = \left(\sum_j D_j^{6}\,\frac{t_{dj}}{t_{mj}}\right)^{1/6}
\qquad \text{(Fórmula 5)}
\end{equation}

para a dose diária. A tensão $m_z D_{zd}$ (Anexo C, Fórmula C.2) entra no
fator de risco $R$ (Fórmula C.3) do mesmo jeito que $S_d^A$ entra em
$R^A$. O ciclo de trabalho de 40\% afeta $D_z$ pelo mesmo motivo que afeta
$S$: picos ausentes do intervalo não gravado não entram na soma.

Isso importa porque o software aceita arquivo \code{.m} de campo também na
aba Cláusula 5/Anexo C (\code{annexc.assess\_from\_segments}), calculando
$D_z$ e $R$ literais sobre o dado de campo. O resultado é numericamente
válido só se o trecho gravado (40\% do tempo) contiver os picos que
definem o dia -- premissa não verificável com o hardware atual. A
Cláusula 5 resolve o problema de sincronismo do dado de campo; não resolve
o problema do ciclo de trabalho.

\subsection{Ausência de sincronismo entre os três módulos}
\label{sec:sync}

Avila et al. (2026), Seção 3, relatam que cada microcontrolador usa seu
próprio relógio interno, o que causa deriva temporal entre os três sinais
(assento, encosto, pé). Citam Elson e Römer (2003) sobre sincronização em
redes de sensores distribuídos e propõem uma arquitetura mestre-escravo
com sinal de sincronismo compartilhado, ainda não implementada.

O Anexo A exige os pontos de contato como entrada simultânea de um único
sistema linear: as funções de transferência combinam a contribuição de
cada ponto no mesmo instante. Sem referência de tempo compartilhada entre
as três unidades, o arquivo do assento e o do encosto do mesmo dia podem
estar defasados por uma quantidade desconhecida, crescente ao longo do
dia. O software trata cada posição como registro independente e não funde
assento, encosto e pé num evento multiaxial.

\subsection{Perda de dados por falha mecânica}
\label{sec:loss}

Avila et al. (2026), Seção 3, atribuem a maior parte da perda de dado em
campo (usuários de cadeira de rodas, 24 h contínuas de rotina real) a
ruptura de fiação: o protótipo inicial usava cabos finos e conectores de
rede adaptados, sensíveis a esforço mecânico repetido e a chuva/umidade.
Em algumas sessões só um dos três acelerômetros ficou ativo por parte do
período.

Arquivos de um mesmo dia podem ter lacunas irregulares além do ciclo de
trabalho de 40\%, e as três posições podem ter durações e períodos
cobertos diferentes entre si. O pipeline de campo
(\code{fielddata.load\_day}, \code{FieldDay.span\_hours}) soma a duração
por arquivo em vez de assumir um período contínuo.

\subsection{Calibração do sensor}
\label{sec:calib}

O MPU6050 é um acelerômetro MEMS de baixo custo, não um instrumento de
referência. Duas fontes de incerteza já registradas no código:

\begin{itemize}
\item A curva de correção de resposta em frequência (``CONEM2024'', em
\code{fielddata.CONEM\_SLOPE}/\code{CONEM\_INTERCEPT}) foi levantada
contra um acelerômetro de referência ICP 353B34 usando uma única unidade
do kit. A variação entre unidades não foi medida.
\item O desvio entre o módulo de gravidade medido parado e o valor de
referência local (\unit{9{,}784}{m/s^2} em Belo Horizonte) pode ser erro
de ganho ou de offset; com uma única orientação, os dois não se
distinguem. O MPU6050 admite offset de até \unit{80}{mg} no eixo z
(cerca de \unit{0{,}78}{m/s^2}), suficiente para 8\% de desvio no módulo.
Por isso essa correção vem desligada por padrão
(\code{fielddata.calibrate\_from\_gravity}); separar as duas causas exige
calibração em seis posições (cada eixo para cima e para baixo).
\end{itemize}

\subsection{A norma foi calibrada para outra população}
\label{sec:populacao}

As funções de transferência do Anexo A (Hinz et al., 2008) e a tabela de
resistência óssea do Anexo C (Fórmula C.4) vêm de população sem
deficiência, majoritariamente motoristas profissionais europeus. O kit de
Avila et al. (2026) mede usuários de cadeira de rodas, população com
osteoporose sublesional documentada na literatura clínica. Estender o
Anexo A a esse público exigiria tabela de resistência óssea e função de
transferência levantadas nessa população, não emprestadas de motoristas.
O único trabalho nessa direção é o modelo biodinâmico cinemático de Brown
et al. (2017), com sujeitos de lesão medular real, ainda incompatível com
o formato de dose/risco do Anexo C.

\section{Mudanças propostas}

\subsection{Hardware/firmware}
\begin{itemize}
\item Arquitetura mestre-escravo com sinal de sincronismo compartilhado,
proposta por Avila et al. (2026), não implementada.
\item Cabeamento robusto (parcialmente feito) e placas de circuito
impresso dedicadas, para reduzir conexões soltas.
\item O ciclo de trabalho de 40\% vem do tempo de escrita no cartão SD.
Um cartão ou protocolo de escrita mais rápido, ou um buffer que escreva
o bloco anterior enquanto grava o seguinte, pode reduzir essa perda.
\end{itemize}

\subsection{Calibração}
\begin{itemize}
\item Repetir a calibração de resposta em frequência em cada unidade do
kit, não só na unidade original, e registrar a variação entre unidades.
\item Calibração em seis posições por unidade, para separar ganho de
offset.
\end{itemize}

\subsection{Software}
\begin{itemize}
\item Manter o caminho de campo (NHO-09/Diretiva/Cláusula 5, bloco a
bloco, sem colagem) como via principal para dado do registrador.
\item Deixar explícito na interface que os três arquivos de um mesmo dia
não são temporalmente comparáveis ponto a ponto.
\item Avisar na aba Cláusula 5/Anexo C, quando o dado de entrada vier de
arquivo \code{.m} de campo, que $D_z$/$R$ herdam o problema da
Seção~\ref{sec:clause5}.
\end{itemize}

\subsection{Metodológica}
\begin{itemize}
\item Função de transferência espinhal específica de população com lesão
medular, nos moldes do que Brown et al. (2017) fizeram para cinemática.
\item Com sincronismo de hardware, reavaliar se vale reconstituir um
Anexo-A-like para o kit de campo, ou se NHO-09/Diretiva bastam.
\end{itemize}

\section{Conclusão}

Os pontos acima vêm do dado de origem: ciclo de trabalho parcial, falta
de sincronismo entre unidades, perda de dado por falha mecânica,
calibração de uma unidade só, população fora da premissa da norma.
Nenhum se resolve só no software. Este apontamento reúne essas causas com
a fonte de cada uma, para que a próxima decisão de hardware, calibração
ou software parta do motivo real.

\section*{Referências}
\begin{itemize}
\item Avila, J.P.R.; Lima II, E.J.; Donadon, L.V.; Duarte, M.L.M.; Augusto,
S.O.; Huang, F.; Pereira, P.T.S. \emph{Development of a Low-Cost
Alternative for Measuring Human Whole-Body Vibration (WBV): Problems and
Perspectives}. CONEM2026-0596, 13\textsuperscript{th} Brazilian Congress
of Mechanical Engineering, 2026.
\item Brown, K.; Flashner, H.; McNitt-Gray, J.; Requejo, P. \emph{Modeling
Wheelchair-Users Undergoing Vibrations}. Journal of Biomechanical
Engineering, 139(9), 2017.
\item Elson, J.; Römer, K. \emph{Wireless Sensor Networks: A New Regime
for Time Synchronization}. Computer Communication Review, 2003.
\item Goeking, V.L.V. \emph{Estudo de um sistema de aquisição de dados
portátil para a avaliação de vibrações de corpo inteiro em usuários de
cadeiras de rodas}. Trabalho de conclusão de curso, Departamento de
Engenharia Mecânica, UFMG, 2023.
\item Goeking, V.L.V.; Duarte, M.L.M.; Lima II, E.J.; Donadon, L.V.
\emph{Projeto, desenvolvimento e análise de dispositivo portátil de
baixo custo para medição de exposição à vibração de corpo inteiro em
usuários de cadeiras de rodas}. Anais do 12\textsuperscript{o} Congresso
Brasileiro de Engenharia Mecânica, 2024.
\item Hinz, B.; Seidel, H.; Hofmann, J.; Menzel, G. \emph{The significance
of using anthropometric parameters and postures of European drivers as a
database for finite-element models when calculating spinal forces during
whole-body vibration exposure}. International Journal of Industrial
Ergonomics, 38, 816--843, 2008.
\item ISO 2631-1:1997. \emph{Mechanical vibration and shock -- Evaluation
of human exposure to whole-body vibration -- Part 1: General
requirements}. International Organization for Standardization, Geneva.
\item ISO 2631-5:2018. \emph{Mechanical vibration and shock -- Evaluation
of human exposure to whole-body vibration -- Part 5: Method for
evaluation of vibration containing multiple shocks}. International
Organization for Standardization, Geneva.
\end{itemize}

\end{document}
